\section{ARCH-15-ORDINARY-TARGET: a published-only target for the ordinary case}
\label{arch15-ordinary-target}

\paragraph{Scope.}
This supplementary reduction does not construct the unknown mixed
cross-cycle for an arbitrary ordinary K3 surface or cubic fourfold.
It does remove the all-CM outside preprint from the construction and
full realization of a suitable \emph{chosen} target over a finite
extension. The resulting mixed Tate criterion for that target is
necessary and sufficient without \(\mathrm{AB}\). Comparison with an
arbitrary previously fixed target remains subject to the scope
distinction in ARCH-15-INPUT-AUDIT. Inputs: published Honda--Tate,
semisimplicity and ordinary Tate for the varieties at issue, and
Milne (2001), Appendix A.7, second Example on author p.~21; the
remaining algebraic arguments have tag none. A, M and all three named
outside CM preprints are unused.

\subsection{A Frobenius packet and its rational projector}

Let \(q=p^a\). Suppose a nonzero semisimple rational Frobenius packet
has an integral weight-two characteristic polynomial \(f\), with
exactly one root of normalized \(p\)-adic slope 0, exactly one of
slope 2, and all remaining roots of slope 1, and has no roots
\(q\zeta\) with \(\zeta\) a root of unity. Here the slopes include
multiplicities and are measured relative to \(v_p(q)=a\), using any
fixed embedding into \(\overline{\mathbf Q}_p\).

\textbf{Lemma.} The polynomial \(f\) is irreducible and occurs with
multiplicity one. If \(\alpha\) is one root and
\(K=\mathbf Q(\alpha)\), then \(K\) is CM,
\(\mathbf Q(\alpha^n)=K\) for every \(n>0\), and there is a rational
algebraic projector on the degree-\((1,1)\) part of
\(h(B\times B)\), for an abelian variety \(B/\mathbf F_{q^2}\), whose
image has Frobenius roots precisely the conjugates of \(\alpha^2\),
each once. The projector is asserted in the homological category and
has a rational Chow representative. Every power of \(B\) satisfies
the published Tate conjecture.

\textbf{Proof.} Let \(f_0\) be an irreducible factor containing the
unique slope-0 root. Since \(\bar\alpha=q^2/\alpha\), the conjugate
slope-2 root belongs to that same factor. Its multiplicity is one
because the total slope-0 multiplicity is one. Every root of any
other factor would have slope 1 at every place above \(p\). Dividing
it by \(q\) gives an algebraic integer with all complex absolute
values one: outside \(p\), division by \(q\) changes no integrality,
and at \(p\) the asserted slopes ensure integrality. Kronecker's
theorem would make that quotient a root of unity, contrary to the
hypothesis. Thus \(f=f_0\). Complex conjugation on \(K\) is the
involution \(\alpha\mapsto q^2/\alpha\); it has no real embedding
in this case and its fixed field is totally real. Hence \(K\) is CM.

For any \(n>0\), put \(L=\mathbf Q(\alpha^n)\). Choose the embedding
of \(K\) giving slope zero. Every embedding of \(K\) extending its
restriction to \(L\) sends \(\alpha\) to a root with the same nth
power and thus also slope zero. Uniqueness of the slope-zero root
forces \([K:L]=1\). The same count shows that the corresponding
place \(v_0\mid p\) has local degree one. Its conjugate also has
local degree one; all other places have slope 1 for \(\alpha\) relative
to \(q\).

Now regard \(\alpha\) as an integral weight-one Weil number for the
field of size \(q^2\). Honda--Tate supplies a simple abelian variety
\(B/\mathbf F_{q^2}\) with this label. The slopes of its label are
\(0,1/2,1\); its center is \(K\). Every local invariant of its
division algebra is either 0 or 1/2 modulo \(\mathbf Z\), so its
index \(d\) is 1 or 2. The preceding stabilization of \(K\), and the
unchanged slope ratios under finite extension, show that the index and
the simple dimension remain unchanged after every finite extension.
Thus \(B\) is geometrically simple. The geometric endomorphism algebra
is already defined over a finite extension and has these same center,
local invariants and index. Milne's Appendix A.7 example applies:
there is a degree-one prime of center at slope 0, its conjugate at
slope 1, and all other primes have slope 1/2. It proves
\(P(B)=L(B)\), hence Tate for every power of \(B\).

Here is the rational projector, including its multiplicity check.
Let \(U=h^1(B)\), and let \(D\) be its endomorphism division algebra;
depending on the covariance convention it is the opposite of
\(\operatorname{End}^0(B)\), which has the same index and whose Brauer
class also has order at most two. Endomorphisms of \(B\) induce its
action by rational algebraic correspondences. The degree-1 projector
itself can, for the present homological purpose, be taken as a
rational polynomial in the multiplication-by-\(m\) correspondence:
the distinct eigenvalues \(m^i\), \(0\leq i\leq2\dim B\), isolate
degree 1. Thus no unproved Chow lifting of a numerical projector is
needed.

The finite etale algebra \(K\otimes_{\mathbf Q}K\) has a distinguished
factor \(K\), defined by multiplication \(x\otimes y\mapsto xy\).
Let \(e_\Delta\) be the corresponding central idempotent. After a
splitting field its components select exactly pairs of identical
embeddings \((\sigma,\sigma)\). Its action on \(U\otimes U\) is
algebraic. On the selected component the acting central simple algebra
is
\[
 e_\Delta(D\otimes_{\mathbf Q}D)=D\otimes_KD.
 \tag{ARCH-15.O1}
\]
Its Brauer class is \(2[D]=0\), and its dimension over \(K\) is
\(d^4\); hence it is \(\operatorname{Mat}_{d^2}(K)\). Choose a
primitive idempotent \(j\) in this matrix algebra. It too is a
rational algebraic correspondence on the selected tensor component.
Over an algebraic closure of the coefficient field, \(R_\ell U\)
has dimension \(d\) for every embedding \(\sigma:K\to
\overline{\mathbf Q}_\ell\), with Frobenius \(\sigma(\alpha)\).
The diagonal tensor block has dimension \(d^2\) for each \(\sigma\),
and \(\operatorname{Mat}_{d^2}\) acts on it by its standard
representation. Consequently \(j\) has rank one in every such block.
Its image
\[
 N=(U\otimes U,e_\Delta j)
 \tag{ARCH-15.O2}
\]
has dimension \([K:\mathbf Q]\), with eigenvalues
\(\sigma(\alpha)^2\), each once. The rational Chow representative
acts with this same property for every \(\ell\ne p\). This proves
the lemma. \(\square\)

\subsection{Application to ordinary K3 surfaces and cubic fourfolds}

Let \(X\) first be an ordinary K3 surface. Extend its finite field so
that all geometric Neron--Severi classes are defined. Semisimplicity
and the published divisor Tate theorem identify the complement
\(T_\ell(X)\) with the packet having no root-of-unity eigenvalues after
twist by (1). The ordinary Newton polygon has slopes
\(0,1^{20},2\), so that packet has exactly the slope pattern in the
lemma. Its rational characteristic polynomial exists: remove the
finite-order algebraic factor from the rational \(H^2\) polynomial.
The lemma proves the irreducibility and multiplicity assertions here
directly, without silently treating a sum of rational factors as one
simple packet. After the further quadratic extension,
\(R_\ell N\simeq T_\ell(X)\). Together with the explicit
Neron--Severi, degree-zero and degree-four Tate parts this gives a
chosen rational homological target for all of \(h(X)\).

For an ordinary cubic fourfold, first split its algebraic middle
cohomology over a finite extension. Published Tate in codimension 2
and semisimplicity make the root-of-unity part of \(H^4(2)\) exactly
that algebraic part, with a nondegenerate intersection pairing. Choose
actual cycle representatives and its inverse Gram matrix to split it
off. This gives a rational Chow projector over that chosen extension;
no canonical Chow descent of this middle projector is asserted.
If \(\gamma\) is a remaining \(H^4\) eigenvalue, put
\(\alpha=\gamma/q\). The ordinary \(H^4\) slopes are
\(1,2^{21},3\), so \(\alpha\) is integral and its packet has slopes
\(0,1,2\) with both extreme multiplicities one. Apply the lemma.
Over the quadratic extension, \(N\otimes\mathbb L\) has eigenvalues
\(q^2\alpha^2=\gamma^2\), and is the desired transcendental
weight-four target. All remaining cohomology is a sum of Tate motives.

The target in either case lies in the tensor category generated by
\(B\), with Tate twists: a polarization and its intersection degree
split \(\mathbb L\) from \(h^2(B)\). Thus it has a finite untwisted
presentation by powers of \(B\); no unrelated auxiliary elliptic
factor, and no unproved mixed-product theorem for such a factor, is
necessary. Tate on every power of \(B\), plus semisimplicity and
duality, gives full homological realization for this target category.
For this chosen presentation the argument of CC-XC03 therefore gives,
\emph{using only published inputs},
\[
 \begin{gathered}
 \text{forward rational cross-cycles to the chosen target exist}\\
 \Longleftrightarrow\quad
 \text{the projected mixed Tate map is onto}\\
 \Longleftrightarrow\quad
 \text{Tate in codimension \(\dim X\) holds on its finite family
 of products \(X\times B^{n_\nu}\)}.
 \end{gathered}
 \tag{ARCH-15.O3}
\]
The projected domain and codimension are exactly those in
ARCH-15-CUBIC-GATE or ARCH-15-K3-GATE after the stated base extension.
An isomorphism into this target implies Tate on the unprojected mixed
products because both factors then lie in the category generated by
\(B\). Conversely their Tate maps project to the required Hom space.
Rational coefficient selection and simultaneous certification use the
same restricted-category argument as CC-XC12--15.

If the cycles in (ARCH-15.O3) are supplied over the extension, induction
and trace give a forward certificate for the original \(X/k\) into
an induced abelian--Artin target. Its abelian restrictions of scalars
are geometrically powers of \(B\): field conjugates of an abelian
variety over a finite field have the same point counts, hence the same
Frobenius polynomial and are isogenous by Tate's theorem. Full Tate
within that induced sector follows from powers of \(B\) after finite
extension and trace averaging. This is a sufficient certificate over
the original field; the lemma does not manufacture the missing mixed
cycle. It identifies an additional place where a blanket outside
all-CM hypothesis can be replaced by a published theorem.
